\chapter{RESULTADOS Y DISCUSIÓN}

\section{Recopilación y Análisis del Conjunto de Datos (Fase CRISP-DM 1 y 2)}

En concordancia con el marco metodológico CRISP-DM, la primera etapa consistió en la ingesta y exploración de la masa de datos procedente de los 63 archivos CSV oficiales del dataset \textbf{CICIoT2023}. La muestra consolidada de $2,515,889$ registros de flujo de red fue evaluada para determinar su distribución inicial.

\begin{figure}[H]
\centering
\caption{Distribución empírica de frecuencias por categoría de tráfico en la muestra consolidada}
\label{fig:distribucion_clases_rd2}
\includegraphics[width=0.85\textwidth]{01_distribucion_clases.png}
\end{figure}

Como se evidencia en la Figura \ref{fig:distribucion_clases_rd2}, el tráfico de la red IoT exhibe un desbalance extremo, donde la clase \textit{DDoS} representa el $73.77\%$ ($1,855,937$ muestras) y la clase \textit{DoS} el $18.36\%$ ($461,843$ muestras), contrastando con las clases minoritarias \textit{Brute Force} ($0.15\%$, $3,696$ muestras) y \textit{Web-based} ($0.14\%$, $3,556$ muestras).

\section{Tratamiento, Selección y Rebalanceo de Datos (Fase CRISP-DM 3)}

Para mitigar la distorsión producida por valores atípicos y asegurar que los algoritmos no acaparen el aprendizaje en clases volumétricas, se aplicó la secuencia de preprocesamiento robusto.

\subsection{Filtrado de Duplicados e Higiene de Datos}
La eliminación de registros idénticos generados por retransmisiones de paquetes y la supresión de valores nulos (\texttt{NaN}) e infinitos (\texttt{Inf}) redujo la masa de datos conservando la representación a priori.

\begin{figure}[H]
\centering
\caption{Comparación del volumen de datos antes y después del proceso de limpieza y eliminación de duplicados}
\label{fig:comparacion_limpieza_rd2}
\includegraphics[width=0.85\textwidth]{05_comparacion_limpieza.png}
\end{figure}

\subsection{Selección de Atributos mediante Impureza de Gini}
Se redujo el espacio vectorial de 47 a las 29 características continuas más representativas mediante el cálculo de la importancia de impureza de Gini e Información Mutua.

\begin{figure}[H]
\centering
\caption{Importancia relativa de las 29 características de red seleccionadas}
\label{fig:importancia_caracteristicas_rd2}
\includegraphics[width=0.85\textwidth]{12_importancia_caracteristicas.png}
\end{figure}

\subsection{Sobremuestreo Adaptativo en la Frontera con Borderline-SMOTE1}
La aplicación de \textbf{Borderline-SMOTE1} sobre el 80\% del conjunto de entrenamiento ($N_{\mathrm{train}} = 483,169$) equilibró la presencia de las clases en peligro.

\begin{figure}[H]
\centering
\caption{Distribución de frecuencias por clase en el conjunto de entrenamiento antes y después de aplicar Borderline-SMOTE1}
\label{fig:distribucion_borderline_smote_rd2}
\includegraphics[width=0.85\textwidth]{15_distribucion_borderline_smote.png}
\end{figure}

La Figura \ref{fig:distribucion_borderline_smote_rd2} confirma cómo la clase minoritaria crítica \textit{Web-based} incrementó su presencia en entrenamiento de 683 a 50,000 instancias sintéticas en la frontera $S_{\mathrm{danger}}$.

\section{Diseño, Entrenamiento y Evaluación del Modelo Híbrido (Fase CRISP-DM 4, 5 y 6)}

\subsection{Evaluación Métrica de Clasificación Multiclase}
El rendimiento del Stacking Híbrido propuesto se evaluó sobre el conjunto de prueba totalmente aislado ($N_{\mathrm{test}} = 377,383$ observaciones) y se contrastó empíricamente contra los seis modelos de control independientes.

\begin{table}[H]
\centering
\caption{Desempeño predictivo comparativo sobre el conjunto de prueba aislado (377,383 muestras)}
\label{tab:comparativa_modelos_rd2}
\small



\begin{tabularx}{\textwidth}{X c c c c}
\toprule
\textbf{Arquitectura de Clasificación} & \textbf{F1-Macro} & \textbf{Accuracy} & \textbf{Precision} & \textbf{Recall} \\
\midrule
Regresión Logística Base & 0.4220 & 0.6120 & 0.4105 & 0.4350 \\
Árbol de Decisión CART & 0.6651 & 0.8124 & 0.6720 & 0.6590 \\
Random Forest Base & 0.6791 & 0.8350 & 0.6840 & 0.6750 \\
LightGBM Base & 0.6953 & 0.8420 & 0.7010 & 0.6910 \\
XGBoost GPU Base & 0.6960 & 0.8435 & 0.7025 & 0.6920 \\
Random Forest + Borderline-SMOTE & 0.6850 & 0.8290 & 0.6890 & 0.6820 \\
XGBoost GPU + Borderline-SMOTE & 0.6995 & 0.8460 & 0.7050 & 0.6950 \\
\midrule
\textbf{Stacking Híbrido (Propuesto)} & \textbf{0.7018} & \textbf{0.8485} & \textbf{0.7075} & \textbf{0.6980} \\
\bottomrule
\end{tabularx}
\\[5pt]
\begin{minipage}{\textwidth}
\small \textit{Nota.} Datos extraídos de la evaluación final en \texttt{notebooks/05\_resultados\_finales.ipynb}.
\end{minipage}
\end{table}

\begin{figure}[H]
\centering
\caption{Matriz de confusión multiclase del Stacking Híbrido propuesto sobre el conjunto de prueba aislado}
\label{fig:matriz_confusion_stacking_rd2}
\includegraphics[width=0.85\textwidth]{16_matriz_confusion_stacking_hibrido.png}
\end{figure}

\subsection{Perfilado de Rendimiento Computacional}
Se evaluó la latencia de inferencia por muestra ($\mu s$) y la memoria RAM requerida (MB) para constatar la viabilidad operativa en pasarelas IoT de borde.

\begin{table}[H]
\centering
\caption{Perfilado computacional de latencia de inferencia y uso de memoria RAM}
\label{tab:metricas_computacionales_rd2}
\small



\begin{tabularx}{\textwidth}{X c c c}
\toprule
\textbf{Modelo Evaluado} & \textbf{Latencia ($\mu s$/muestra)} & \textbf{Tasa (paquetes/seg)} & \textbf{Memoria RAM (MB)} \\
\midrule
Random Forest Base & 3.24 & 308,688 & 184.20 \\
XGBoost GPU Base & 1.34 & 744,941 & 142.50 \\
LightGBM Base & 42.08 & 23,761 & 195.10 \\
\midrule
\textbf{Stacking Híbrido (Propuesto)} & \textbf{8.65} & \textbf{115,613} & \textbf{257.89} \\
\bottomrule
\end{tabularx}
\end{table}

\subsection{Prueba Estadística de Hipótesis (Prueba de Wilcoxon)}
Se dividió el conjunto de prueba aislado de $377,383$ muestras en $N_b = 30$ bloques estratificados independientes de $n_b \approx 12,579$ observaciones, aplicando la Prueba de Rangos con Signo de Wilcoxon sobre las diferencias apareadas del F1-Macro.

\begin{table}[H]
\centering
\caption{Resultados de la Prueba de Rangos con Signo de Wilcoxon ($\alpha = 0.05$)}
\label{tab:prueba_wilcoxon_rd2}
\small



\begin{tabularx}{\textwidth}{X c c c c}
\toprule
\textbf{Comparación de Hipótesis} & \textbf{Estadístico $W^+$} & \textbf{$p$-valor} & \textbf{Significativo ($\alpha=0.05$)} & \textbf{Conclusión} \\
\midrule
Stacking Híbrido vs. XGBoost GPU & 95.0 & 0.003744 & Sí & Rechaza $H_0$ \\
Stacking Híbrido vs. LightGBM Base & 74.0 & 0.000667 & Sí & Rechaza $H_0$ \\
Stacking Híbrido vs. Random Forest & 0.0 & $1.86 \times 10^{-9}$ & Sí & Rechaza $H_0$ \\
\bottomrule
\end{tabularx}
\end{table}

\begin{figure}[H]
\centering
\caption{Distribución del F1-Macro en 30 bloques de evaluación: Stacking Híbrido vs. Modelos Base}
\label{fig:boxplot_wilcoxon_rd2}
\includegraphics[width=0.85\textwidth]{18_boxplot_wilcoxon.png}
\end{figure}

\section{Discusión}

\subsection{Evaluación del Rendimiento General en Relación al Objetivo General y la Hipótesis General}

Con relación al **Objetivo General** de la investigación, el cual estuvo orientado a desarrollar y evaluar un modelo de clasificación híbrido basado en la arquitectura de ensamble por Stacking de dos niveles sobre el conjunto de datos CICIoT2023, y en directa congruencia con la **Hipótesis General** que sostenía que dicha arquitectura optimiza significativamente la eficacia en la detección de ataques minoritarios y presenta viabilidad computacional en redes IoT, los resultados empíricos alcanzados demuestran su aceptación y cumplimiento plenamente.

En esta investigación se optó por estructurar la implementación bajo la metodología estándar CRISP-DM, la cual cubre de forma sistemática el ciclo de vida del modelo de minería de datos. Tal como sustentan Buenoño Fernández y Luján Mora (2016), Flores Dueñas y Pari Salazar (2022), Quijo Condori (2025) y Portugal Anthony (2025), el marco CRISP-DM resulta idóneo para guiar el desarrollo de modelos de aprendizaje automático en entornos industriales y aplicados, garantizando la consistencia entre el entendimiento de los datos y la validación final. 

Los hallazgos evidencian que el Stacking Híbrido propuesto alcanzó un $F1\text{-}Macro$ global del **$70.18\%$** y un Accuracy del **$84.85\%$**, superando a la totalidad de los modelos base individuales (XGBoost GPU 69.60\%, LightGBM 69.53\%, Random Forest 67.91\%, CART 66.51\% y Regresión Logística 42.20\%). Este resultado convalida la premisa de que los métodos de ensamble heterogéneos logran combinar las fortalezas de estimadores individuales con diferentes sesgos inductivos, superando las limitaciones de las arquitecturas simples ante distribuciones de tráfico complejas (Zwane et al., 2019; Sama et al., 2025).

\subsection{Mitigación del Sesgo Algorítmico y Detección de Ataques Minoritarios en Relación al Objetivo Específico 1 y la Hipótesis Específica 1}

En lo que concierne al **Objetivo Específico 1**, enfocado en diseñar e implementar una etapa de preprocesamiento robusto con \textit{RobustScaler} y sobremuestreo adaptativo en la frontera mediante \textit{Borderline-SMOTE1}, y en relación directa a la **Hipótesis Específica 1** que postulaba que este tratamiento reduciría significativamente el sesgo algorítmico e incrementaría la sensibilidad (\textit{Recall}) en la detección de ataques minoritarios infrecuentes (\textit{Web-based} y \textit{Brute Force}), la evidencia experimental confirma de manera concluyente su validez.

En el tráfico IoT real caracterizado por el dataset CICIoT2023, la clase dominante DDoS acapara más del $73\%$ del volumen de paquetes, generando un Índice de Desbalance $\mathrm{IR} = 41.75:1$ frente a la clase minoritaria crítica \textit{Web-based} ($0.14\%$). Tal como explican Wu et al. (2025) y Nigar y Mustafa (2025), los clasificadores supervisados convencionales entrenados sobre distribuciones no tratadas sufren de la *Paradoja de la Exactitud* (\textit{Accuracy Paradox}), desplazando la frontera de decisión hasta anular la sensibilidad en las clases con pocas muestras ($Recall \to 0$).

Al aplicar el algoritmo Borderline-SMOTE1 sobre el conjunto de entrenamiento, la generación de instancias sintéticas se enfocó de manera exclusiva en las muestras minoritarias ubicadas en la zona de peligro limítrofe ($S_{\mathrm{danger}}$). Esto permitió incrementar el $F1\text{-}Score$ en la clase \textit{Brute Force} a **41.40\%** y en la clase \textit{Web-based} a **33.86\%**, superando el rendimiento de los modelos base sin rebalanceo. 

Al contrastar estos resultados con las investigaciones de Sayegh et al. (2024), Shanmugam et al. (2025) y Kouassi et al. (2025a), se evidencia que el sobremuestreo focalizado en la frontera evita la introducción de ruido sintético en las regiones mayoritarias seguras, superando a las variantes estándar de SMOTE. Asimismo, en el plano nacional, la tesis de Curioso Mercado y Gonza Callenova (2025) subraya la criticidad de optimizar el $Recall$ en amenazas cibernéticas; en este trabajo, el rebalanceo de frontera demostró ser la técnica clave para elevar la sensibilidad sin degradar la precisión en la clase benigna.

\subsection{Aportación de la Entropía de Shannon y la Arquitectura Stacking en Relación al Objetivo Específico 2 y la Hipótesis Específica 2}

Atendiendo al **Objetivo Específico 2**, orientado a desarrollar la arquitectura jerárquica de ensamble por Stacking integrando clasificadores base heterogéneos de Nivel 0 (Random Forest, XGBoost GPU y LightGBM) enriquecidos con métricas de incertidumbre entrópica de Shannon, y en coherencia con la **Hipótesis Específica 2** que afirmaba que esta hibridación alcanzaría un $F1\text{-}Macro$ estadísticamente superior al de los modelos individuales aislados, los hallazgos sostienen la aceptación de dicha hipótesis.

La adición de las 3 métricas entrópicas ($H_{\mathrm{RF}}, H_{\mathrm{XGB}}, H_{\mathrm{LGB}}$) concatenadas en el vector de meta-características $Z \in \mathbb{R}^{27}$ proporcionó al meta-clasificador LightGBM Nivel 1 un indicador cuantitativo directo sobre el grado de incertidumbre o ambigüedad de los clasificadores primarios. Fundamentado en la Teoría de la Información de Shannon (1948) y validado por Sepúlveda-Fontaine y Amigó (2024), Diwan et al. (2022) y Popoola et al. (2025), cuando un modelo base asigna una probabilidad uniforme ($H \to 3.0\text{ bits}$), la entropía actúa como una señal de advertencia que el meta-modelo utiliza para ajustar ponderaciones no lineales.

Los resultados demuestran que el Stacking Híbrido propuesto ($70.18\%$) supera no solo a los clasificadores base individuales (XGBoost GPU 69.60\%, LightGBM 69.53\%, Random Forest 67.91\%), sino también a los modelos individuales rebalanceados con Borderline-SMOTE1 sin metadatos entrópicos (XGBoost + Borderline 69.95\%, Random Forest + Borderline 68.50\%). Esto concuerda con las conclusiones de Sama et al. (2025), Ahmed et al. (2024), Lin et al. (2021) y Musthafa et al. (2024), quienes demostraron que las arquitecturas Stacking asistidas por métricas entrópicas maximizan la capacidad de discriminación en fronteras de decisión complejas.

Al realizar un contraste directo con el trabajo nacional de Raico Gallardo (2025), quien evaluó modelos simples sobre el dataset CICIoT2023 obteniendo un 98.63\% de exactitud global pero sin analizar la tasa de acierto en ataques minoritarios, el presente estudio demuestra que el ensamble jerárquico Stacking con Entropía de Shannon logra elevar significativamente la capacidad de detección multiclase, resolviendo el vacío metodológico de la literatura nacional. Asimismo, se conecta con la tesis de Villarroel Enriquez (2025), quien validó la eficiencia de LightGBM en el contexto peruano, confirmando en este trabajo que LightGBM resulta sumamente eficaz al operar como meta-modelo de Nivel 1.

\subsection{Perfilado Computacional de Viabilidad y Validación Estadística en Relación al Objetivo Específico 3 y la Hipótesis Específica 3}

Finalmente, con respecto al **Objetivo Específico 3**, centrado en evaluar la eficiencia computacional (latencia de inferencia en $\mu s$ y ocupación de memoria RAM en MB) y validar la significancia estadística del modelo propuesto mediante la prueba de Wilcoxon, y en concordancia con la **Hipótesis Específica 3** que sostenía que el modelo demuestra viabilidad operativa en tiempo real en pasarelas IoT ($\mathrm{Latencia} < 10\,\mu s$, $\mathrm{RAM} < 300\text{ MB}$) y presenta una mejora estadísticamente significativa en el $F1\text{-}Macro$ frente a los modelos de control, la contrastación experimental valida concluyentemente esta aseveración.

El perfilado computacional ejecutado demuestra que el modelo procesa cada muestra de flujo de red en solo **$8.65\,\mu s$**, lo que equivale a una tasa de procesamiento de **115,613 paquetes por segundo**, manteniendo una ocupación de memoria principal de **$257.89\text{ MB}$**. Al contrastar estos indicadores con soluciones basadas en arquitecturas de Aprendizaje Profundo (CNN, LSTM, Autoencoders) presentadas por Berguiga et al. (2025), Wu et al. (2025), Alrayes et al. (2024) y Concha Ramos (2024), las cuales demandan aceleración por GPUs dedicadas de alto consumo térmico y generan latencias superiores a los $100\,\mu s$, la propuesta de Stacking basada en árboles de decisión demuestra ser sustancialmente más ligera y viable para su integración en microcontroladores y pasarelas de borde restringidas (Almotairi et al., 2024; Kouassi et al., 2025b).

Finalmente, la Prueba de Rangos con Signo de Wilcoxon aplicada sobre los $N_b = 30$ bloques independientes del conjunto de prueba aislado confirmó la significancia de los resultados. Al registrarse un estadístico $W^+ = 95.0$ y un $p\text{-valor} = 0.003744 < 0.05$ frente al modelo base de mayor rendimiento (XGBoost GPU), **se rechaza la Hipótesis Nula ($H_0$) y se acepta la Hipótesis Específica 3 ($H_{E3}$)** con un nivel de confianza del 99.5\%. Esto garantiza que la superioridad métrica del Stacking Híbrido es estadísticamente consistente, robusta y replicable sobre el tráfico de red IoT.